\subsection{Hyperbolic Evolution Problems}\label{sec:hyperbolic-evolution-problems}

\subsubsection{Abstract Second-Order Evolution Problems}\label{subsec:abstract-second-order-evolution-problems}

Let $H$ be a separable real Hilbert space and let $V$ be another separable Hilbert space, which is continuously and densely imbedded in $H$. Moreover, let $A\in C^1([0,T],\mathcal{L}(V,V^*))$, and let $a(t,u,v)=-(A(t)u,v)$ be the associated quadratic form. We assume that $a$ is symmetric:
\begin{equation}\label{eq:abstract-hyperbolic-form-symmetry}
a(t,u,v)=a(t,v,u),
\end{equation}
and that \eqref{eq:abstract-parabolic-form-coercivity} holds; i.e., there are positive constants $a$ and $b$ such that
\begin{equation}\label{eq:abstract-hyperbolic-form-coercivity}
a(t,u,u)\geq a\|u\|_V^2-b\|u\|_H^2.
\end{equation}
We consider the evolution problem
\begin{equation}\label{eq:abstract-hyperbolic-evolution-initial-value-problem}
\ddot u(t)=A(t)u+f(t),\ u(0)=u_0,\ \dot u(0)=u_1.
\end{equation}
Our goal is the following result.

\begin{Theorem}
\label{thm:abstract-hyperbolic-existence-uniqueness}
Assume that $f\in L^1((0,T),H)$, $u_0\in V$, $u_1\in H$ and that $A$ is as described above. Then there exists a unique weak solution $u\in C([0,T],V)\cap C^1([0,T],H)$ of the evolution problem \eqref{eq:abstract-hyperbolic-evolution-initial-value-problem}.
\end{Theorem}

For the proof, we shall need to proceed in several steps. We shall first give a proof of existence. This proof is similar in spirit to the one given in the previous section for parabolic problems. We first derive an energy equation which yields an a priori estimate for a solution. Then we use an approximation scheme and uniform energy estimates to obtain a solution as a weak limit. The energy estimate yields uniqueness only if more regularity of the solution is assumed; we shall therefore need a separate argument to prove uniqueness. Finally, the existence proof we give will only show that the solution is in $L^\infty((0,T),V)\cap W^{1,\infty}((0,T),H)$, and a separate argument is needed to establish continuity.\par We remark that if $u$ has the regularity guaranteed by the theorem, it does not follow that $f\in L^1((0,T),H)$. Hyperbolic problems differ from parabolic problems in the fact that they lack coercive properties; solutions are not ``as smooth as the data will allow.\textquotesingle\textquotesingle\ This makes these problems in many respects harder than parabolic equations.

\subsubsection{Existence of a Solution}\label{subsec:hyperbolic-solution-existence}

We begin with a formal energy estimate. In \eqref{eq:abstract-hyperbolic-evolution-initial-value-problem}, we take the inner product with $\dot u$. (This is of course not justified; even if we already knew that the theorem is true, $\dot u(t)$ would not be in $V$!). This yields
\begin{equation}\label{eq:abstract-hyperbolic-formal-energy-pairing}
(\ddot u,\dot u)+a(t,u,\dot u)=(f,\dot u).
\end{equation}
We integrate from 0 to $t$, and obtain
\begin{equation}\label{eq:abstract-hyperbolic-formal-energy-identity}
\begin{aligned}a(t,u(t),u(t))+\|\dot u(t)\|_H^2&=a(0,u_0,u_0)+\|u_1\|_H^2\\&\quad+\int_0^t a^{\prime}(s,u(s),u(s))+2(f(s),\dot u(s))\,\mathrm{d}s.\end{aligned}
\end{equation}
Here we have set
\begin{equation}\label{eq:abstract-hyperbolic-time-derivative-quadratic-form}
a^{\prime}(t,u,v)=\dfrac{\partial}{\partial t}a(t,u,v)=-(\dot A(t)u,v).
\end{equation}
From \eqref{eq:abstract-hyperbolic-formal-energy-identity}, one easily derives an estimate of the form
\begin{equation}\label{eq:abstract-hyperbolic-solution-apriori-bound}
\|u(t)\|_V+\|\dot u(t)\|_H\leq C(\|u_0\|_V+\|u_1\|_H+\|f\|_{L^1((0,T),H)}).
\end{equation}
We leave the details of this argument as an exercise (Problem~\ref{ex:hyperbolic-energy-apriori-estimate}).

Next, we construct a sequence of approximate problems for which an analogue of \eqref{eq:abstract-hyperbolic-formal-energy-identity} holds. Let $\{\phi_n\mid n\in\mathbb{N}\}$ be a set of linearly independent vectors in $V$ such that the span of the $\phi_n$ is dense. Let $V_n$ be the span of $\phi_1,\phi_2,\ldots,\phi_n$, let $P_n$ be the orthogonal projection from $H$ onto $V_n$ and let $\Pi_n$ be the orthogonal projection from $V$ onto $V_n$. We now seek
\begin{equation}\label{eq:abstract-hyperbolic-galerkin-finite-span-solution}
u_n(t)=\sum_{i=1}^n\alpha_i(t)\phi_i
\end{equation}
satisfying
\begin{equation}\label{eq:abstract-hyperbolic-galerkin-initial-value-system}
\begin{gathered}(\ddot u_n,\phi_i)=(A(t)u_n,\phi_i)+(f(t),\phi_i),\ i=1,\ldots,n,\\u_n(0)=\Pi_nu_0,\\\dot u_n(0)=P_nu_1.\end{gathered}
\end{equation}
This is a system of second-order ODEs for the $\alpha_i(t)$, which has a unique solution. The energy equation \eqref{eq:abstract-hyperbolic-formal-energy-identity} holds for $u_n$ (with the same derivation as above), and hence one obtains uniform bounds for $u_n\in L^\infty((0,T),V)$ and $\dot u_n\in L^\infty((0,T),H)$. We can extract a weakly-* convergent subsequence, which has a limit $u\in L^\infty((0,T),V)\cap W^{1,\infty}((0,T),H)$. For simplicity, we shall again use the notation $u_n$ to denote the weakly-* convergent subsequence.

It remains to be shown that $u$ is a solution. Let
\begin{equation}\label{eq:abstract-hyperbolic-terminal-zero-test-space}
X=\{\psi\in C^1([0,T],V)\mid\psi(T)=0\}.
\end{equation}
Functions of the form $\psi=\sum_{i=1}^n\alpha_i(t)\phi_i$, where $n\in\mathbb{N}$, are dense in $X$. If $\psi$ is any function of this form, we obtain, for $m\geq n$,
\begin{equation}\label{eq:abstract-hyperbolic-galerkin-test-identity}
\int_0^T(\ddot u_m,\psi)\,\mathrm{d}t=\int_0^T(A(t)u_m,\psi)+(f(t),\psi(t))\,\mathrm{d}t,
\end{equation}
which yields after an integration by parts
\begin{equation}\label{eq:abstract-hyperbolic-galerkin-integrated-test-identity}
\int_0^T-(\dot u_m,\dot\psi)+a(t,u_m,\psi)\,\mathrm{d}t=\int_0^T(f,\psi)\,\mathrm{d}t+(\dot u_m(0),\psi(0)).
\end{equation}
Here we can take the limit, which yields
\begin{equation}\label{eq:abstract-hyperbolic-limit-prescribed-initial-velocity}
\int_0^T-(\dot u,\dot\psi)+a(t,u,\psi)\,\mathrm{d}t=\int_0^T(f,\psi)\,\mathrm{d}t+(u_1,\psi(0)).
\end{equation}
By density, this identity actually holds for every $\psi\in X$. If we restrict $\psi$ to test functions, it follows that the differential equation in \eqref{eq:abstract-hyperbolic-evolution-initial-value-problem} holds in the sense of $V^*$-valued distributions.\par We now note that weak-* convergence in $L^\infty((0,T),H)$ implies weak convergence in $L^2((0,T),H)$. Thus, it follows that $u_n$ converges weakly in $H^1((0,T),H)$. By the Sobolev imbedding theorem, the mapping $u\to u(0)$ is continuous from $H^1((0,T),H)$ to $H$; hence $u_n(0)$ converges weakly in $H$ to $u(0)$. Since, on the other hand, $u_n(0)$ converges strongly in $V$ to $u_0$, it follows that $u(0)=u_0$.

It follows from the differential equation that $\ddot u\in L^1((0,T),V^*)$, so that $\dot u\in C([0,T],V^*)$ (cf. Problem~\ref{ex:wone-one-continuous-closed-interval}). Hence $\dot u(0)$ is meaningful as an element of $V^*$. Repeating the derivation which led to \eqref{eq:abstract-hyperbolic-galerkin-integrated-test-identity} (with $u$ in place of $u_m$), we find
\begin{equation}\label{eq:abstract-hyperbolic-limit-solution-initial-velocity}
\int_0^T-(\dot u,\dot\psi)+a(t,u,\psi)\,\mathrm{d}t=\int_0^T(f,\psi)\,\mathrm{d}t+(\dot u(0),\psi(0)).
\end{equation}
By comparing with \eqref{eq:abstract-hyperbolic-limit-prescribed-initial-velocity}, we conclude that $\dot u(0)=u_1$.

\subsubsection{Uniqueness of the Solution}\label{subsec:hyperbolic-solution-uniqueness}

It would be easy to conclude uniqueness from \eqref{eq:abstract-hyperbolic-formal-energy-identity}. The problem is, however, that the derivation of \eqref{eq:abstract-hyperbolic-formal-energy-identity} requires more smoothness of the solution than we have. We can circumvent this difficulty by deriving an energy equation for time-integrated quantities. Let $u$ be a solution of \eqref{eq:abstract-hyperbolic-evolution-initial-value-problem} for $f=0$, $u_0=u_1=0$. Fix $s\in(0,T)$ and let
\begin{equation}\label{eq:abstract-hyperbolic-uniqueness-integrated-test-function}
v(t)=\begin{cases}-\int_t^s u(r)\,\mathrm{d}r,&t<s,\\0,&t\geq s.\end{cases}
\end{equation}
We multiply \eqref{eq:abstract-hyperbolic-evolution-initial-value-problem} by $v$ and integrate. This yields
\begin{equation}\label{eq:abstract-hyperbolic-uniqueness-integrated-zero-identity}
0=\int_0^T(\ddot u-A(t)u,v)\,\mathrm{d}t=\int_0^T a(t,u,v)-(\dot u,\dot v)\,\mathrm{d}t;
\end{equation}
note that this integration by parts is permissible and the boundary terms vanish since $\dot u(0)=v(T)=0$. Using the definition of $v$, we conclude
\begin{equation}\label{eq:abstract-hyperbolic-uniqueness-primitive-energy-pairing}
\int_0^s a(t,v,\dot v)-(\dot u,u)\,\mathrm{d}t=0.
\end{equation}
Carrying out the integration, we find
\begin{equation}\label{eq:abstract-hyperbolic-uniqueness-primitive-energy-identity}
a(0,v(0),v(0))+\|u(s)\|_H^2=-\int_0^s a^{\prime}(t,v,v)\,\mathrm{d}t.
\end{equation}
We conclude an estimate of the form
\begin{equation}\label{eq:abstract-hyperbolic-uniqueness-primitive-energy-bound}
\|v(0)\|_V^2+\|u(s)\|_H^2\leq C\left(\int_0^s\|v(t)\|_V^2\,\mathrm{d}t+\|v(0)\|_H^2\right).
\end{equation}
Setting $w(t)=\int_0^t u(r)\,\mathrm{d}r=v(t)-v(0)$, we conclude
\begin{equation}\label{eq:abstract-hyperbolic-uniqueness-zero-initial-primitive-bound}
\|w(s)\|_V^2+\|u(s)\|_H^2\leq C\left(\int_0^s\|w(t)-w(s)\|_V^2\,\mathrm{d}t+\|w(s)\|_H^2\right).
\end{equation}
We now use the inequalities
\begin{equation}\label{eq:abstract-hyperbolic-primitive-difference-triangle-bound}
\|w(t)-w(s)\|_V^2\leq2(\|w(t)\|_V^2+\|w(s)\|_V^2)
\end{equation}
and
\begin{equation}\label{eq:abstract-hyperbolic-primitive-cauchy-schwarz-bound}
\|w(s)\|_H^2\leq s\int_0^s\|u(t)\|_H^2\,\mathrm{d}t.
\end{equation}
(The latter follows from the definition of $w$ and Cauchy-Schwarz.) By using these two inequalities in \eqref{eq:abstract-hyperbolic-uniqueness-zero-initial-primitive-bound}, we find
\begin{equation}\label{eq:abstract-hyperbolic-uniqueness-stepping-gronwall-bound}
(1-2Cs)\|w(s)\|_V^2+\|u(s)\|_H^2\leq K\int_0^s\|w(t)\|_V^2+\|u(t)\|_H^2\,\mathrm{d}t
\end{equation}
with some new constant $K$. From this, one easily concludes that $w=u=0$ as long as $s<1/(2C)$ (cf. Gronwall\textquotesingle s inequality). Since the constant $C$ is independent of the starting time, we can use a stepping argument to conclude that $w=u=0$ everywhere in $[0,T]$.

\subsubsection{Continuity of the Solution}\label{subsec:hyperbolic-solution-continuity}

We already know that $u\in C([0,T],H)$ and $\dot u\in C([0,T],V^*)$. The following lemma allows us to draw a further conclusion from this:

\begin{Lemma}
\label{lem:hilbert-triple-weak-continuity}
Let $V$ and $H$ be Hilbert spaces such that $V$ is continuously and densely imbedded in $H$. Assume that $u\in L^\infty((0,T),V)\cap C([0,T],H)$. Then $u(t)\in V$ for every $t\in[0,T]$ and $u(t)$ is weakly continuous; i.e., $(f,u(t))$ is a continuous function of $t$ for every $f\in V^*$.
\end{Lemma}

\begin{Proof}
We shall establish that for each $t\in(0,T)$
\begin{equation}\label{eq:hilbert-triple-weak-continuity-pointwise-bound}
\|u(t)\|_V\leq\|u\|_{L^\infty((0,T),V)}.
\end{equation}
Suppose not. Since $H$ is dense in $V^*$, there exists $f\in H$ such that
\begin{equation}\label{eq:hilbert-triple-weak-continuity-separating-functional}
(f,u(t))>\|f\|_{V^*}\|u\|_{L^\infty((0,T),V)}.
\end{equation}
Since $u\in C([0,T],H)$, we conclude that
\begin{equation}\label{eq:hilbert-triple-weak-continuity-neighborhood-separation}
(f,u(s))>\|f\|_{V^*}\|u\|_{L^\infty((0,T),V)}
\end{equation}
for $s$ in some neighborhood of $t$, say for $|s-t|<\epsilon$. Define now $g(s)=f$ for $|s-t|<\epsilon$ and $g(s)=0$ otherwise. Then, we find
\begin{equation}\label{eq:hilbert-triple-weak-continuity-holder-contradiction}
\int_0^T(g(s),u(s))\,\mathrm{d}s>\|g\|_{L^1((0,T),V^*)}\|u\|_{L^\infty((0,T),V)}.
\end{equation}
This is a contradiction of Hölder\textquotesingle s inequality. Hence $u(t)$ is a bounded function taking values in $V$.\par Consider now $f\in V^*$. Then there exists a sequence $f_n\in H$ such that $f_n\to f$ in $V^*$. It follows that $(f_n,u(t))$ converges uniformly to $(f,u(t))$. Since $(f_n,u(t))$ is continuous, $(f,u(t))$ is continuous.
\end{Proof}

Using the lemma, we conclude that the solution $u$ of \eqref{eq:abstract-parabolic-form-coercivity} is weakly continuous with values in $V$ and $\dot u$ is weakly continuous with values in $H$.

Let us now recall the construction of $u$ in Section~\ref{subsec:hyperbolic-solution-existence}. The solution $u$ was the limit of a sequence $u_n$, and for each $u_n$, we have
\begin{equation}\label{eq:abstract-hyperbolic-galerkin-solution-energy-identity}
\begin{aligned}a(t,u_n(t),u_n(t))+\|\dot u_n(t)\|_H^2&=a(0,\Pi_nu_0,\Pi_nu_0)+\|P_nu_1\|_H^2\\&\quad+\int_0^t a^{\prime}(s,u_n(s),u_n(s))+2(f(s),\dot u_n(s))\,\mathrm{d}s.\end{aligned}
\end{equation}
Consider now any fixed $s\in(0,T]$. The quantity
\begin{equation}\label{eq:abstract-hyperbolic-sup-energy-equivalent-norm}
\sup_{t\in[0,s]}[a(t,u,u)+\|\dot u(t)\|_H^2]
\end{equation}
is equivalent to the square of the norm in $L^\infty((0,s),V)\times L^\infty((0,s),H)$. Since balls are weak-*-compact and hence weak-*-closed, we conclude from \eqref{eq:abstract-hyperbolic-galerkin-solution-energy-identity} that, in the limit $n\to\infty$:
\begin{equation}\label{eq:abstract-hyperbolic-limit-sup-energy-upper-bound}
\begin{aligned}\sup_{t\in[0,s]}a(t,u(t),u(t))+\|\dot u(t)\|_H^2&\leq a(0,u_0,u_0)+\|u_1\|_H^2\\&\quad+\limsup_{n\to\infty}\int_0^s[a^{\prime}(s,u_n(s),u_n(s))+2(f(s),\dot u_n(s))]\,\mathrm{d}s.\end{aligned}
\end{equation}
By letting $s$ tend to zero, we find
\begin{equation}\label{eq:abstract-hyperbolic-initial-right-limsup-energy}
\limsup_{t\to0+}a(t,u(t),u(t))+\|\dot u(t)\|_H^2\leq a(0,u_0,u_0)+\|u_1\|_H^2.
\end{equation}
Since, on the other hand, $u(t)\in V$ and $\dot u(t)\in H$ are weakly continuous, we have
\begin{equation}\label{eq:abstract-hyperbolic-initial-right-liminf-energy}
\liminf_{t\to0+}a(t,u(t),u(t))+\|\dot u(t)\|_H^2\geq a(0,u_0,u_0)+\|u_1\|_H^2.
\end{equation}
It follows that the quantity $a(t,u(t),u(t))+\|\dot u(t)\|_H^2$ is right continuous at $t=0$. However, we might just as well have taken any other time as the initial time; hence we have right continuity everywhere. Finally, since our equation is invariant under time reversal, all our results apply to the time-reversed problem as well. Hence $a(t,u(t),u(t))+\|\dot u(t)\|_H^2$ is also left continuous.

We now note that
\begin{equation}\label{eq:abstract-hyperbolic-strong-continuity-energy-difference}
\begin{aligned}
&a(t,u(t)-u(s),u(t)-u(s))+\|\dot u(t)-\dot u(s)\|_H^2\\
&\quad=a(t,u(t),u(t))-2a(t,u(t),u(s))+a(s,u(s),u(s))\\
&\qquad+a(t,u(s),u(s))-a(s,u(s),u(s))\\
&\qquad+(\dot u(t),\dot u(t))-2(\dot u(t),\dot u(s))+(\dot u(s),\dot u(s)).
\end{aligned}
\end{equation}
By exploiting the weak continuity of $u$ and $\dot u$ and the continuity of $a(t,u,u)+\|\dot u\|^2$, we see that the right-hand side of \eqref{eq:abstract-hyperbolic-strong-continuity-energy-difference} tends to zero as $s\to t$. Hence the left-hand side also tends to zero; this implies continuity of $u(t)\in V$ and $\dot u(t)\in H$.

\Needspace{5\baselineskip}\textbf{Problems}

\begin{enumerate}[label=\thesection.\arabic*.,ref=\thesection.\arabic*,start=6,series=book-problems]

\item\label{ex:hyperbolic-energy-apriori-estimate} Use \eqref{eq:abstract-hyperbolic-formal-energy-identity} to derive \eqref{eq:abstract-hyperbolic-solution-apriori-bound}.

\item\label{ex:hyperbolic-higher-regularity-discussion} Discuss higher regularity of solutions in analogy to Section~\ref{subsec:parabolic-solution-regularity}.

\item\label{ex:hyperbolic-bounded-lower-order-perturbation} Establish an existence and uniqueness result for the following perturbation of \eqref{eq:abstract-hyperbolic-evolution-initial-value-problem}: $\ddot u=A(t)u+B(t)u$, where $A$ is as above and $B$ is a bounded operator from $V$ to $H$.

\item\label{ex:hyperbolic-second-order-pde-applications} Apply the result of this section to the second-order analogues of the examples discussed in Section~\ref{subsec:parabolic-pde-applications}.

\item\label{ex:wone-one-continuous-closed-interval} Let $I$ be an open interval. Show that $W^{1,1}(I)\subset C(\overline I)$.

\item\label{ex:hyperbolic-entire-galerkin-sequence-convergence} In our construction of the solution, we took a subsequence of $u_n$. Show that in fact the whole sequence $u_n$ converges.

\end{enumerate}
